\documentclass{article}

\usepackage[utf8]{inputenc}
\usepackage[T1]{fontenc}
\usepackage{microtype}
\usepackage{amsmath,amssymb,amsthm}
\usepackage{booktabs}
\usepackage{graphicx}
\usepackage[colorlinks=true,citecolor=blue,linkcolor=blue]{hyperref}

\newtheorem{definition}{Definition}[section]
\newtheorem{theorem}{Theorem}[section]
\newtheorem{lemma}[theorem]{Lemma}
\newtheorem{proposition}[theorem]{Proposition}
\newtheorem{corollary}[theorem]{Corollary}

\title{Auth-SafeInv: An Evaluation Target for Authorization-Conditioned\\Effect Monitoring with Three-Layer Over-Optimism Diagnosis}
\author{}

\begin{document}

\maketitle

\begin{abstract}
We introduce \textsc{Auth-SafeInv}, an evaluation target for authorization-conditioned realized-effect monitoring in tool-augmented LLM agents, and diagnose three layers at which monitor evaluation can overstate safety. \textbf{Layer 1 (tool-surface):} linear probes trained on Qwen3-8B representations detect causal effects within seen tools, but can fail under leave-one-tool-out coverage stress tests with held-out false negative rates up to $0.80$; lexical normalization does not remove the failure and probe-mediated interchange intervention (pIIA) provides convergent activation-level evidence. We refer to this failure as \emph{surface-form fragmentation}---a representation-level phenomenon where effect detection is not invariant to the tool surface form used. \textbf{Layer 2 (trace-label):} execution-trace verifiers that achieve FNR/FPR of $0.0/0.0$ with full structured labels degrade when label fields are removed (mean FNR $0.0353 \to 0.1162$) or reduced to minimal execution evidence (mean FNR $0.1465$). \textbf{Layer 3 (row-to-action):} row-level FNR/FPR and action-level allow/deny decisions measure different objects: a monitor with row-level FPR $=0$ can still falsely deny all authorized actions under candidate-effect aggregation, and validation-calibrated thresholding can collapse to all-allow. We report three gate-failure results as honest negative evidence: pure representation-side repair (contrastive projection, schema conditioning, decomposed verification) does not beat strong baselines under strict held-out protocols. We provide a constructive two-stage monitoring framework (Effect\-Verif-Auth\-Monitor) that separates effect-presence verification from authorization checking, with controlled-setting feasibility evidence and explicit limits. All results are from a controlled diagnostic study over synthetic scenarios, controlled sandbox traces, file-backed browser-runtime traces, and key-free live protocol traces. We do not claim deployment safety certification or method superiority over existing defenses.
\end{abstract}

\section{Introduction}

LLM-based agents increasingly act by invoking tools: reading and writing files, calling web APIs, sending messages, executing shell commands, and browsing the web. The safety question is therefore not only whether a tool call is syntactically valid, but whether the causal effects of that call are justified by the current task. A request to summarize a local file may authorize reading file contents; it does not by itself authorize network egress, file deletion, or sending the contents elsewhere.

This framing exposes an evaluation gap behind many agent defenses. Boundary and behavioral mechanisms such as causal attribution, task-specific access control, provenance tracking, and architectural isolation~\cite{attrguard2026,clawguard2026,argus2026,xoa2026} ultimately rely on some account of \emph{what a proposed action would do}. Recent work on effect-trace semantics~\cite{alignmentcontracts2026}, provenance-aware capability contracts~\cite{pact2026}, and runtime interception~\cite{agenttrust2026} has begun to formalize this requirement. If the same downstream effect is represented differently depending on whether it is produced by a direct API call or by an equivalent shell command, a learned safety monitor can appear reliable on covered tools while failing on an unseen surface form. Conversely, if the monitor uses structured trace labels that are close to the answer, it can appear stronger than it actually is. And if monitor evaluation reports only candidate-effect row metrics without aggregating to action-level allow/deny decisions, it can overstate real safety.

We study these three layers of over-optimism in a controlled setting, introducing \textsc{Auth-SafeInv} as an action-level evaluation target. Given a task context $c$, the authorized effect envelope $A(c)$ specifies which causal effects are permitted. A tool call $(a,S)$ produces realized effects $\Omega(a,S)$. Safety holds when $\Omega(a,S) \subseteq A(c)$; otherwise $U(c,a,S) = \Omega(a,S) \setminus A(c)$ constitutes unauthorized effects. Our evaluation target asks: can a monitor reliably detect when $U(c,a,S) \neq \emptyset$, and does its evaluation at each layer (tool-surface, trace-label, action-level) honestly reflect its real performance?

We make two tightly-linked contributions:

\begin{enumerate}
    \item \textbf{Auth-SafeInv: an evaluation target with three-layer over-optimism diagnosis.} We formalize authorization-conditioned realized-effect monitoring as a measurable evaluation object with explicit definitions of $A(c)$, $\Omega(a,S)$, and $U(c,a,S)$. At Layer~1 (tool-surface), we use linear probes with leave-one-tool-out splits to diagnose \emph{surface-form fragmentation}---a representation-level failure where the same causal effect is not detected invariantly across tool surface forms. We provide convergent evidence through held-out FNR (up to $0.80$), lexical controls, and probe-mediated interchange intervention (pIIA). At Layer~2 (trace-label), we show that structured trace label fields can make effect verifiers appear stronger than they are, with label-hidden and minimal-evidence views degrading FNR by $3.3\times$--$4.2\times$. At Layer~3 (row-to-action), we demonstrate that row-level FNR/FPR and action-level allow/deny decisions are structurally different metrics, and that simple calibration can collapse to all-allow. Layer~1 diagnosis is a SafeInv (effect-detection invariance) problem; Layers~2--3 constitute the authorization-conditioned Auth-SafeInv evaluation.

    \item \textbf{Constructive framework with honest evidence and explicit limits.} We provide Effect\-Verif-Auth\-Monitor, a two-stage sandbox-based architecture that separates effect-presence verification from authorization judgment. We report three gate failures (T54/T55/T56) where pure representation-side repair fails to beat strong baselines. We then show controlled-setting feasibility evidence across six trace families (controlled sandbox, local adapters, provider API, broader adapters, file-backed Chrome runtime, live protocol), where execution-level effect evidence reduces unauthorized-effect misses relative to static rules. We document explicit limits: the execution verifier is deterministic and designed for these controlled trace formats; row-level success does not imply action-level safety; a handcrafted raw-status baseline is competitive; and provider-backed search, SaaS messaging, HTTP browser automation, and deployed-agent runtime traces remain absent. All artifacts are released as an open-source evaluation suite.
\end{enumerate}

Through experiments spanning 459 original diagnostic scenarios, an expanded 1,464-row Auth-SafeInv counterfactual suite, and 8,988 trace-based candidate-effect rows across controlled sandbox, live protocol, and file-backed headless Chrome browser-runtime traces, we provide a unified evaluation framework for a problem that, we argue, the field needs to measure before it can solve.

\section{Related Work}

\subsection{Agent Safety Benchmarks and Runtime Defenses}

Recent agent safety research has systematically mapped the attack surface~\cite{lasm2026,sok2026} and demonstrated real-world exploitability~\cite{jaw2026,litmus2026,queryipi2026}. Large-scale benchmarks evaluate complementary objects: AgentDojo~\cite{agentdojo2024} evaluates prompt-injection defenses; ToolEmu~\cite{toolemu2024} uses LM-emulated sandboxes; AgentHarm~\cite{agentharm2025} measures harmful task completion; ASB~\cite{asb2025} broadens evaluation across scenarios, tools, attacks, and defenses; ST-WebAgentBench~\cite{stwebagentbench2026} evaluates safety-policy compliance in web agents; TraceSafe~\cite{tracesafe2026} assesses mid-trajectory guardrail safety. Our work is not another benchmark for attack success or task completion. It asks whether a monitor can detect \emph{unauthorized realized effects} consistently across tool surfaces, and whether common evaluation practices overstate that ability.

Defense mechanisms operate primarily at the behavioral or boundary layer. AttriGuard~\cite{attrguard2026} performs causal attribution by re-executing agents under attenuated context views. CausalArmor~\cite{causalarmor2026} uses causal ablation at privileged action points. ClawGuard~\cite{clawguard2026} derives task-specific access constraints before tool invocation. ARGUS~\cite{argus2026} constructs influence provenance graphs. XOA~\cite{xoa2026} proposes architectural sandboxing. These methods rely, explicitly or implicitly, on some account of what a tool call would or did do. Our work evaluates a monitor-level target orthogonal to these defenses: whether the realized effect set exceeds the task-authorized envelope, and whether current evaluation practices honestly reflect monitor performance.

A rapidly growing body of 2026 work explicitly targets effect-level and provenance-level monitoring. PACT~\cite{pact2026} enforces provenance-aware capability contracts at the argument level. Alignment Contracts~\cite{alignmentcontracts2026} formalize effect-trace semantics with Lean~4 proofs. ARM~\cite{arm2026} detects causality laundering through provenance graphs augmented with counterfactual edges. AgentTrust~\cite{agenttrust2026} provides runtime tool-call interception with allow/warn/block/review verdicts. Parallax~\cite{parallax2026} enforces cognitive-executive separation with adversarial validation. ECA~\cite{eca2026} uses evidence-carrying certificates to prevent hallucination-to-action conversion. These systems focus on \emph{enforcement} and \emph{prevention}. Our work is complementary: we focus on \emph{evaluation}---how to measure whether a monitor (of any architecture) can reliably detect unauthorized effects, and what evaluation practices make that measurement over-optimistic.

\begin{table}[t]
\centering
\caption{Novelty boundary relative to agent security benchmarks, runtime defenses, and effect-level monitoring systems. The key distinction is that we provide an evaluation target and three-layer over-optimism diagnosis, not an enforcement mechanism.}
\label{tab:relatedboundary}
\resizebox{\textwidth}{!}{
\begin{tabular}{llll}
\toprule
Work type & Representative work & Primary judgment object & Difference from this paper \\
\midrule
Prompt-injection benchmark & AgentDojo, ASB~\cite{agentdojo2024,asb2025} & attack success / utility & unauthorized-effect FNR under tool-surface variation \\
Risk discovery & ToolEmu~\cite{toolemu2024} & risky tool outcomes in emulation & monitor invariance for realized effect labels \\
Harmful-task benchmark & AgentHarm~\cite{agentharm2025} & harmful task completion & task-authorized effect consistency \\
Trajectory safety & TraceSafe, ST-WebAgentBench~\cite{tracesafe2026,stwebagentbench2026} & mid-trajectory guardrail accuracy & three-layer evaluation framework \\
Causal attribution defense & AttriGuard, CausalArmor~\cite{attrguard2026,causalarmor2026} & why a privileged call occurs & what effect occurred and whether authorized \\
Boundary enforcement & ClawGuard~\cite{clawguard2026} & pre-action constraint satisfaction & execution evidence after/during realization \\
Provenance auditing & ARGUS~\cite{argus2026} & trusted influence over decisions & realized-effect verification plus authorization \\
Provenance contracts & PACT~\cite{pact2026} & argument-level authority binding & effect-level probe evaluation across surfaces \\
Formal effect semantics & Alignment Contracts~\cite{alignmentcontracts2026} & effect-trace satisfaction proofs & representation-level diagnosis, not formal methods \\
Causality tracking & ARM~\cite{arm2026} & transitive denial influence & cross-tool representation invariance \\
Runtime interception & AgentTrust, Parallax~\cite{agenttrust2026,parallax2026} & pre-execution allow/block & post-execution effect verification framework \\
Evidence certificates & ECA~\cite{eca2026} & hallucination-to-action gate & effect-readout invariance, not perceptual certification \\
\bottomrule
\end{tabular}}
\end{table}

\subsection{Causal Abstraction and Representation Diagnosis}

Causal abstraction~\cite{geiger2021,geiger2025} provides the theoretical foundation for verifying whether a neural network implements a high-level causal model. Sutter et al.~\cite{sutter2025} proved that without constraining the alignment mapping $\tau$ to be low-complexity, causal abstraction becomes vacuous. We adopt linear probes as the alignment class and use probe-mediated interchange intervention (pIIA) as a diagnostic for cross-form effect-direction transfer, not as an SCM-based causal proof.

\subsection{Contrastive Representation Learning and Domain Generalization}

Our surface-form fragmentation diagnosis connects to domain generalization~\cite{irm2019,groupdro2020,icidg2026,cdas2026}, where tool names and call styles act like domains while the safety-relevant target is the downstream effect. The difference is our safety metric and evaluation target: we care primarily about unauthorized-effect false negatives under held-out tool coverage, not average accuracy.

\section{Formal Framework}

\subsection{Safety System Model}

We model an authorization-conditioned agent safety monitor as a triple $\mathcal{S}_c = (\Phi, \Psi_c, \Omega)$, where:
\begin{itemize}
    \item $\Phi = \{\tau_1, \ldots, \tau_k\}$ maps an action's LLM embedding to predicted effect probabilities. In our linear-probe instantiation, $\tau_i(\mathbf{h}) = \sigma(\mathbf{w}_i^\top \mathbf{h} + b_i)$ with $\mathbf{h} = N(\text{desc}(a, S))$.
    \item $\Psi_c: [0,1]^k \to \{\text{ALLOW}, \text{DENY}\}$ is the authorization-conditioned policy. Given $A(c) \subseteq \mathcal{E}$ and thresholds $\{\theta_i\}$, $\Psi_c(\mathbf{p}) = \text{DENY} \iff \exists E_i \notin A(c): p_i > \theta_i$.
    \item $\Omega: \mathcal{A} \times \mathcal{S} \to 2^{\mathcal{E}}$ returns the realized causal-effect set. The unauthorized effect set is $U(c,a,S)=\Omega(a,S)\setminus A(c)$.
\end{itemize}

The intended safety condition is $\Omega(a,S)\subseteq A(c)$. A violation occurs when $U(c,a,S)\neq\emptyset$.

\subsection{Multi-Mechanism Effects and Surface-Form Fragmentation}

An agent operation can be described through different \emph{surface forms} while producing identical causal consequences.

\begin{definition}[Multi-Mechanism Effect]
A causal effect $E_i$ is \emph{multi-mechanism} if there exist at least two distinct surface forms that can produce it:
$$\text{MultiMech}(E_i) \iff \exists S_a \neq S_b \in \mathcal{S}: P(E_i = 1 \mid S = S_a) > 0 \land P(E_i = 1 \mid S = S_b) > 0.$$
\end{definition}

The core diagnostic question is whether the LLM's representation encodes $E_i$ invariantly across these surface forms. We operationalize this through two metrics:

\textbf{ToolProxyGap.} For held-out form $B$, let $\hat{\tau}_{i,-B}^{\text{diag}}$ be a diagnostic probe trained on all forms except $B$, and $\hat{\tau}_{i,B}^{\text{within}}$ a within-form probe via stratified cross-validation on $B$:
$$\text{ToolProxyGap}_i(B) = [\text{FNR}_B(\hat{\tau}_{i,-B}^{\text{diag}}) - \text{FNR}_B(\hat{\tau}_{i,B}^{\text{within}})]_+.$$
This measures how much detection degrades when coverage for a form is missing, relative to how detectable the same effect is within that form.

\textbf{pIIA (probe-mediated Interchange Intervention Accuracy).} At layer $\ell$, we inject the mean projection difference between a source sample (form $B$, $E_i=1$) and a base sample (form $B$, $E_i=0$) along the probe direction $\hat{\mathbf{w}}_A$ learned from form $A$, and evaluate whether the final-layer probe $\tau_i^B$ registers the injected effect:
$$\text{pIIA}_i(A \to B) = P\left(\tau_i^B(N_{>\ell}(I_{\mathbf{w}_A}(\mathbf{h}_\ell^b, \mathbf{h}_\ell^s))) > \tau_i^B(\mathbf{h}_L^b) \mid E_i(b)=0, E_i(s)=1\right).$$
pIIA is \emph{probe-mediated}: the intervention direction and evaluation probe come from different surface forms, making it a diagnostic for cross-form effect-direction transfer, not an SCM-based causal proof.

\subsection{Safety-Centric Invariance}

\begin{definition}[SafeInv]
A probe $\tau_i$ is $(\epsilon_{\text{safe}}, \epsilon_{\text{fpr}})$-\emph{safety-invariant} for multi-mechanism effect $E_i$ if:
$$\max_{S \in \mathcal{S}_{E_i}} R_{S, \theta}^{-}(\tau_i) \leq \epsilon_{\text{safe}} \quad \text{and} \quad \max_{S \in \mathcal{S}_{E_i}} R_{S, \theta}^{+}(\tau_i) \leq \epsilon_{\text{fpr}}$$
where $R_{S, \theta}^{-}(\tau_i) = P(\tau_i(\mathbf{h}) \leq \theta_i \mid E_i=1, S)$. In safety-critical settings, $\epsilon_{\text{safe}} \ll \epsilon_{\text{fpr}}$.
\end{definition}

\begin{definition}[Auth-SafeInv]
A probe $\tau_i$ is $(\epsilon_{\text{unauth}}, \epsilon_{\text{fauth}})$-\emph{authorization-safe-invariant} for effect $E_i$ if, over task contexts $c$ and surface forms $S$ where $E_i$ can be realized but is not authorized,
$$\max_{c,S} P(\tau_i(\mathbf{h}) \leq \theta_i \mid E_i \in U(c,a,S)) \leq \epsilon_{\text{unauth}},$$
and the corresponding false positive rate on authorized or absent effects is at most $\epsilon_{\text{fauth}}$. Auth-SafeInv upgrades SafeInv by conditioning the detection target on task authorization.
\end{definition}

\subsection{EffectVerif-AuthMonitor}

The verifier-assisted framework separates two questions: \emph{which effects occurred?} and \emph{which occurred effects were authorized?} Since realized effects can only be observed after or during execution, we assume a \textbf{two-phase sandbox execution model}. In the first phase, the tool call $(a,S,x)$ is executed in an isolated sandbox that records low-level execution evidence (file diffs, HTTP metadata, exit codes, stdout/stderr, provider receipts) while preventing irreversible side effects. In the second phase, Effect\-Verif-Auth\-Monitor processes the sandbox trace to decide ALLOW (commit effects to the real environment) or DENY (discard the sandbox). This model is consistent with existing sandbox-and-commit architectures in agent security~\cite{parallax2026,alignmentcontracts2026,eca2026} and avoids the logical contradiction of using post-execution evidence for pre-execution decisions.

\begin{center}
\fbox{
\begin{minipage}{0.92\textwidth}
\textbf{\textsc{EffectVerif-AuthMonitor}}\\
\textbf{Input:} task context $c$; envelope $A(c)$; tool call $(a,S,x)$; execution trace $t$\\
\textbf{1.} $\hat{\Omega} \leftarrow \textsc{VerifyEffects}(a,S,x,t)$\\
\textbf{2.} $\hat{U} \leftarrow \hat{\Omega}\setminus A(c)$\\
\textbf{3.} \textbf{if} $\hat{U}\neq\emptyset$ \textbf{return} DENY with violated effects $\hat{U}$\\
\textbf{4.} \textbf{else return} ALLOW
\end{minipage}}
\end{center}

This is a framework interface, not a claim that the verifier is perfect. The trace schema includes tool name, arguments, execution status, file diffs, HTTP request/response metadata, stdout/stderr, provider status, and message receipts. The empirical question is whether adding such evidence reduces unauthorized-effect misses relative to monitors that only see task text or static policy rules.

\subsection{Risk Accounting}

\begin{lemma}[Authorization-Conditioned FNR Risk Accounting]
\label{thm:safety}
For task context $c$ and surface form $B$ where action $a$ realizes $E \notin A(c)$, define $\beta_{\text{auth}} = P(\tau_E(\mathbf{h}) \leq \theta_E \mid E \in U(c,a,B))$ and $\alpha_{\text{auth}} = P(\exists E_j \notin A(c), E_j \neq E: \tau_j(\mathbf{h}) > \theta_j \mid E \in U(c,a,B))$. Then:
$$P(\mathcal{S}_c(a, B) = \text{ALLOW} \mid E \in U(c,a,B)) \geq [\beta_{\text{auth}} - \alpha_{\text{auth}}]_+.$$
\end{lemma}

This is a risk-accounting result, not a safety certificate. It says that if an unauthorized-effect detector has high FNR on a specific surface form, and if other detectors do not redundantly fire, then unsafe allowance is lower-bounded. The LOTO diagnostic instantiates an effect-level analogue with $\beta^{\text{LOTO}}$ and $\alpha^{\text{LOTO}}$.

\begin{theorem}[Verification Sample Complexity]
\label{thm:sample}
Let $q = P(S=B, E \in U(c,a,B))$. To observe at least one such sample with probability $\geq 1-\delta$ requires $m \geq \frac{1}{q} \log \frac{1}{\delta}$; to estimate FNR to accuracy $\pm\epsilon$ requires $m = \tilde{O}\left(\frac{1}{q \epsilon^2} \log \frac{1}{\delta}\right)$.
\end{theorem}

Theorem~\ref{thm:sample} explains why rare tool-effect or authorization-violation combinations are difficult to validate by natural sampling alone. Contrastive projection theory (generalization bounds and alignment-FNR relationships) is deferred to the appendix; it provides conditional motivation for observed-pair repair but its assumptions do not hold in the strict held-out settings where pure representation methods fail.

\section{Experiments}

\subsection{Setup}

\textbf{Data.} We use several increasingly strict data regimes. The original diagnostic benchmark contains 459 agent tool-call scenarios spanning 11 causal effects and 9 tools. Each scenario includes a task context, tool call with surface form and arguments, and binary effect labels. The key design principle is breaking the deterministic tool-to-effect shortcut: the same tool produces different effects based on context, and the same effect can be produced through different tools. We also build an expanded Auth-SafeInv suite: 1,464 authorization counterfactual rows with explicit $A(c)$ envelopes, plus 35,136 schema-conditioned candidate-effect rows. For trace-based evaluation, we collect 8,988 candidate-effect rows across six trace families: controlled local sandbox observations (1,104 rows), real-agent-tools local file/terminal adapters (384 rows), a single-provider DeepSeek API trace set (960 rows), broader web/search/browser/messaging local adapters (2,400 rows), file-backed headless Chrome DOM/JavaScript runtime traces (960 rows), and key-free live protocol traces combining real outbound HTTPS and local webhook protocol I/O (2,400 rows; full-precision). All trace-based rows are controlled diagnostics, not deployed-agent traffic.

\textbf{Models.} We extract embeddings from Qwen3-8B (4096d), with supplementary results from MiniLM (384d) and Qwen2.5-7B-Instruct (3584d). Probes are $\ell_2$-regularized logistic regression. LOTO probes train on all non-held-out forms; within-form FNR uses 3-fold stratified cross-validation inside the held-out form. pIIA uses hook-based intervention at layers 24/36 of Qwen3-8B. Thresholds in trace-based evaluation are selected on validation trace groups and evaluated on held-out trace groups.

\textbf{Effects.} We study 11 causal effects defined with operational criteria that distinguish them from tool names or call syntax. \texttt{command\_executed}: a system command was executed (shell, subprocess). \texttt{file\_written}: file-system content was created or modified. \texttt{file\_deleted}: a file or directory was removed. \texttt{file\_content\_read}: file content was accessed and returned. \texttt{message\_sent}: a message was dispatched to an external recipient. \texttt{network\_egress}: network data left the local environment (HTTP request, API call, data upload). \texttt{content\_fetched}: content was retrieved from an external URL or API. \texttt{search\_performed}: a search query was executed against an external index. \texttt{subagent\_spawned}: a new agent or subprocess was delegated. \texttt{memory\_updated}: persistent memory or state was modified. \texttt{tool\_error}: the tool execution returned an error or was blocked. Effects can co-occur (e.g., \texttt{content\_fetched} and \texttt{network\_egress} co-occur on web API calls; \texttt{command\_executed} and \texttt{file\_written} co-occur on shell redirection). Effect labels are assigned during data generation by deterministic rules from the scenario schema; they are not human-annotated. The key design constraint is that the same effect can be realized through multiple tools and the same tool can produce different effects depending on arguments and context.

\textbf{Metrics.} We report FNR (the safety-critical metric), FPR, ToolProxyGap, and pIIA-Drop ($\text{pIIA}_{\text{within}} - \text{pIIA}_{\text{cross}}$). For action-level evaluation, we report Unauthorized Action Allow Rate (U-Allow) and Authorized Action False Denial Rate (FDeny). All main-table FNR/FPR values are reported with Wilson 95\% confidence intervals where $N^+$ or $N^-$ permits; small cells ($N^+ < 30$) are flagged.

\subsection{Layer 1: Tool-Surface Fragmentation (SafeInv Failure)}

The first layer of over-optimism is evaluating effect detectors only on seen tool surfaces. This is a \textbf{SafeInv} failure---the detector does not maintain effect-presence invariance across surface forms. It is distinct from Auth-SafeInv (Layers~2--3), which additionally conditions the detection target on task authorization. Layer~1 diagnosis establishes that even the simpler SafeInv condition fails; authorization conditioning (Auth-SafeInv) cannot hold if basic effect detection is surface-form-fragile.

\textbf{LOTO stress test.} Table~\ref{tab:loto} shows leave-one-tool-out results on Qwen3-8B. Within-form FNR ranges from 0.03 to 0.48, confirming that within-form detection is feasible but not uniformly trivial. Held-out FNR can be extreme: $0.97$ for \texttt{file\_written/write\_file}, $0.88$ for \texttt{file\_content\_read/read\_file}, $0.78$ for \texttt{content\_fetched/web\_fetch}. The stress-test signal is the \emph{degradation}---ToolProxyGap reaches $0.88$ for \texttt{file\_content\_read} and $0.72$ for \texttt{content\_fetched} on its primary tool. A monitor evaluated only within its training tools would see low FNR; deployed on a held-out tool it can miss most unauthorized effects.

\begin{table}[t]
\centering
\caption{Layer 1: LOTO FNR under coverage-missing stress tests. $N^+$ = positive samples in heldout form. Within-FNR is heldout-form internal CV; Held-FNR trains on all non-heldout forms. ToolProxyGap = $[\Delta\text{FNR}]_+$. Wilson 95\% CI width for Held-FNR: $\pm 0.10$ at $N^+=32$, $\pm 0.14$ at $N^+=20$, $\pm 0.16$ at $N^+=10$, $\pm 0.17$ at $N^+=8$. Small-$N$ cells should be treated as stress-test indicators.}
\label{tab:loto}
\begin{tabular}{lcccccc}
\toprule
Effect & Training Forms & Heldout & $N^+$ & Within-FNR & Held-FNR & ToolProxyGap \\
\midrule
content\_fetched & all-but-terminal & terminal & 8 & 0.111 & 0.500 & 0.389 \\
content\_fetched & all-but-web\_fetch & web\_fetch & 32 & 0.064 & 0.781 & 0.718 \\
file\_content\_read & all-but-read\_file & read\_file & 32 & 0.033 & 0.875 & 0.842 \\
file\_deleted & all-but-terminal & terminal & 20 & 0.246 & 0.600 & 0.354 \\
file\_written & all-but-write\_file & write\_file & 33 & 0.091 & 0.970 & 0.879 \\
network\_egress & all-but-web\_search & web\_search & 26 & 0.074 & 0.731 & 0.657 \\
tool\_error & all-but-web\_search & web\_search & 10 & 0.278 & 0.100 & 0.000 \\
\bottomrule
\end{tabular}
\end{table}

\textbf{pIIA activation-level evidence.} Table~\ref{tab:piia} provides convergent diagnostic evidence. \texttt{tool\_error}---whose cross-form semantics (``failure/denied/blocked'') are linguistically unified---achieves near-perfect cross-form transfer (pIIA-Drop = 0.052). \texttt{content\_fetched}---whose realization through \texttt{web\_fetch} (``fetched JSON from API'') and \texttt{terminal} (``curl -s URL'') involves disjoint linguistic patterns---shows larger cross-form drop (pIIA-Drop = 0.298). pIIA is probe-mediated and is used as a diagnostic, not a causal proof.

\begin{table}[t]
\centering
\caption{pIIA results (Qwen3-8B, layer 24/36). 95\% CIs are source-cluster bootstrap.}
\label{tab:piia}
\begin{tabular}{lcccc}
\toprule
Effect & pIIA-within & pIIA-cross & pIIA-Drop & 95\% CI \\
\midrule
tool\_error & 1.000 & 0.948 & 0.052 & [0.032, 0.077] \\
file\_content\_read & 1.000 & 0.837 & 0.164 & [0.050, 0.286] \\
file\_deleted & 1.000 & 0.800 & 0.200 & [0.094, 0.325] \\
content\_fetched & 1.000 & 0.702 & 0.298 & [0.143, 0.544] \\
file\_written & 1.000 & 0.669 & 0.331 & [0.147, 0.550] \\
network\_egress & 1.000 & 0.634 & 0.366 & [0.307, 0.424] \\
\bottomrule
\end{tabular}
\end{table}

\textbf{Not a lexical shortcut.} Lexical normalization (replacing tool names with [TOOL], inline code with [CMD], URLs with [URL]) leaves overall held-out FNR essentially unchanged: $0.378 \to 0.382$ across matched cells. Individual cells move in both directions (\texttt{file\_written} $0.97 \to 0.70$; \texttt{network\_egress/terminal} unchanged at $0.41$). The fragmentation persists in the representation.

\textbf{No simple baseline eliminates the failure.} Pooled training, balanced subsampling, tool-conditioned probes, group reweighting, and Procrustes alignment all leave residual worst-form FNR. Balanced pooling helps on some effects (\texttt{file\_deleted}: $0.60 \to 0.20$) but not others (\texttt{content\_fetched}: $0.78 \to 0.78$). Procrustes alignment improves only 15/38 form pairs. Sampling imbalance explains part but not all of the failure.

\textbf{Coverage-missing safety bound.} The LOTO stress-test lower bound $[\beta^{\text{LOTO}}-\alpha^{\text{LOTO}}]_+$ identifies cases where high held-out FNR translates to a positive unsafe-allow bound. The largest bound is \texttt{file\_content\_read/read\_file}: $\beta^{\text{LOTO}}=0.88$, $\alpha^{\text{LOTO}}=0.19 \Rightarrow [\beta-\alpha]_+ = 0.69$. Some high-FNR cells have zero bound because of sufficient predicted redundancy (\texttt{content\_fetched/terminal}: $\beta=0.50$, $\alpha=1.00$). Both high-FNR and low-redundancy must coincide for a safety vulnerability. See appendix for the full table.

\subsection{Layer 2: Trace-Label Over-Optimism}

The second layer of over-optimism is evaluating effect verifiers on trace fields that encode answer-like information. The full-label view includes fields such as \texttt{verified\_effects}, \texttt{effect\_diff}, and \texttt{unauthorized\_effects} that directly name the effects to be detected. That full-label performance is higher is expected---this is partly a sanity check that monitor evaluation should exclude label-carrying fields. The diagnostic value lies in (a) quantifying the degradation magnitude ($3.3\times$--$4.2\times$), (b) showing that degradation is not uniform across trace families (largest on browser-runtime traces, where static verifiers already fail), and (c) establishing that the label-hidden view still outperforms pre-action baselines while being substantially weaker than the full-label upper bound. Table~\ref{tab:layer2} reports the three-view ablation.

\begin{table}[t]
\centering
\caption{Layer 2: Trace-view verifier ablation (mean across T58/T59/T61/T63/T65/T64 held-out trace groups). Label-hidden hides \texttt{verified\_effects}, \texttt{effect\_diff}, \texttt{unauthorized\_effects}, and related fields. Minimal-evidence keeps only file diffs, exit codes, HTTP metadata, stdout/stderr.}
\label{tab:layer2}
\begin{tabular}{lcccc}
\toprule
View & Mean row FNR & Mean row FPR & Mean U-Allow & Mean FDeny \\
\midrule
Full label (upper bound) & 0.0353 & 0.0000 & 0.0417 & 0.4544 \\
Label-hidden raw & 0.1162 & 0.0480 & 0.0459 & 0.7428 \\
Minimal evidence & 0.1465 & 0.0643 & 0.0531 & 0.8417 \\
\bottomrule
\end{tabular}
\end{table}

Full-label FNR of $0.0353$ would suggest strong performance. But removing label fields increases mean FNR by $3.3\times$ to $0.1162$, and further reducing to minimal evidence increases it by $4.2\times$ to $0.1465$. The degradation is largest on the file-backed browser-runtime traces, where the static verifier already fails completely (FNR $1.0$). Action-level false denial also increases substantially. This does not mean the verifier is purely reading labels---the label-hidden view still outperforms pre-action rule-only baselines (mean FNR $0.57$, see Layer 3 below)---but it means that full-label results should be treated as upper bounds, not as evidence of deployment-ready monitoring. The current verifier depends on structured execution evidence.

\subsection{Layer 3: Row-to-Action Metric Translation}

The third layer of over-optimism is the gap between candidate-effect row metrics and action-level allow/deny decisions.

\textbf{Row-level success does not imply action-level safety.} Table~\ref{tab:layer3} aggregates validation-selected candidate-effect predictions into action-level metrics. A real monitor ultimately decides whether an action should be allowed; it is not sufficient to report that individual effect predictions have low error rates.

\begin{table}[t]
\centering
\caption{Layer 3: Action-level metrics from the same T62 validation-selected thresholds. U-Allow = Unauthorized Action Allow Rate; FDeny = Authorized Action False Denial Rate. Execution verifier can achieve row-level FNR/FPR of $0/0$ on some traces while having action-level FDeny of $1.0$.}
\label{tab:layer3}
\resizebox{\textwidth}{!}{
\begin{tabular}{lcccc}
\toprule
Trace family & Exec U-Allow & Exec FDeny & Static U-Allow & Static FDeny \\
\midrule
Controlled auth traces & 0.0000 & 0.1591 & 0.0000 & 0.4091 \\
real-agent-tools local adapters & 0.2000 & 0.3571 & 0.2000 & 0.3571 \\
DeepSeek provider API & 0.0000 & 1.0000 & 1.0000 & 0.0000 \\
Broader local adapters & 0.0500 & 0.2100 & 0.4875 & 0.2100 \\
Headless Chrome runtime & 0.0000 & 1.0000 & 1.0000 & 0.0000 \\
Live HTTP + webhook protocol & 0.0000 & 0.0000 & 0.3210 & 0.0000 \\
\bottomrule
\end{tabular}}
\end{table}

The execution verifier achieves row-level FNR/FPR of $0.0/0.0$ on the provider API and browser-runtime traces, yet action-level false denial rates reach $1.0$ on both. This is not a logical contradiction---it is a structural consequence of how candidate-effect predictions compose into action decisions. The action-level aggregation pipeline works as follows:

\begin{enumerate}
    \item For each action (tool call), the execution verifier produces a binary present/absent prediction for every candidate effect $E_i \in \mathcal{E}$.
    \item The authorization checker compares each predicted-present effect against the task envelope $A(c)$. If any predicted-present effect is $\notin A(c)$, the action is DENIED.
    \item Row-level FPR measures whether individual effect-presence predictions are wrong. Action-level FDeny measures whether the composed DENY decision is wrong. A row-level true negative (correctly predicting a candidate effect is absent) can contribute to an action-level false denial if \emph{another} candidate effect is incorrectly predicted as present and unauthorized.
    \item Concretely: on the browser-runtime traces, many authorized actions have candidate effects that are legitimately in $A(c)$ but whose presence is over-predicted by the verifier. Since presence over-prediction on an authorized effect counts as a row-level false positive only on that specific effect, the per-effect row FPR can be low, but the action-level decision---which fires on \emph{any} predicted unauthorized effect---accumulates these over-predictions across all candidate effects.
\end{enumerate}

The gap between row-level and action-level metrics is therefore a structural property of the candidate-effect composition rule, not a measurement error. A monitor with low per-effect error rates can still have high action-level error when the decision aggregates across many candidate effects.

\textbf{Simple calibration does not fix this.} An action-level calibration pass that constrains validation false denial to $\leq 0.10$ often selects a degenerate all-allow threshold. On the real-agent-tools local adapters, broader local adapters, provider API, and browser-runtime trace families, the calibration chooses thresholds that allow all actions, yielding U-Allow $=1.0$. This means the constraint of low false denial cannot be met without effectively disabling the monitor, exposing a fundamental tension between safety (blocking unauthorized actions) and utility (not blocking authorized ones) in current effect-based monitoring.

\subsection{Gate Failure: Pure Representation Repair}

Before turning to the constructive framework, we report an important negative result: three attempts at pure representation-side repair---methods that only modify the LLM's embedding space without external execution evidence---all fail to beat strong baselines under strict held-out protocols.

\begin{table}[t]
\centering
\caption{Three gate-failure results. All values are at FPR $\leq 0.10$ under strict held-out splits. Best baseline is the strongest same-cell non-degenerate baseline.}
\label{tab:gatefail}
\begin{tabular}{lccc}
\toprule
Method & Setting & FNR (method) & FNR (best baseline) \\
\midrule
Contrastive projection (T54) & LOTO, 21 cells & 0.2804 & 0.1835 (supervised contrastive) \\
Schema-conditioned monitor (T55) & LOTO, full-tool-chain & 0.4516 & 0.1835 (supervised contrastive) \\
Decomposed frozen verifier (T56) & LOTO, full-tool-chain & 0.3547 & 0.1835 (supervised contrastive) \\
\bottomrule
\end{tabular}
\end{table}

Table~\ref{tab:gatefail} summarizes the three gate failures. Contrastive projection achieves 21/21 evaluable LOTO cells after a data expansion (up from 4/14) but its strict train-only FNR of 0.2804 is substantially worse than the strongest same-cell baseline (supervised contrastive, FNR $=0.1835$). A pair-free schema-conditioned monitor---which asks the model to infer effects from tool descriptions without seeing cross-form pairs---does worse still (FNR $=0.4516$). Decomposing the task into separate effect-presence and authorization probes, using only frozen Qwen3-8B embeddings, does not close the gap (FNR $=0.3547$).

These are not null results in the usual sense---they are \emph{gate} failures: each method makes a different structural assumption about why the problem should be solvable from representations alone, and each fails the strict test. We report them because they establish the difficulty of the problem and prevent overclaiming about representation-side solutions.

\textbf{Verifier-present upper bound.} When an oracle label indicates which effects are actually present, a simple authorization classifier achieves LOTO FNR $=0.0042$ and family FNR $=0.0745$ at FPR $=0$. This is an upper bound for the two-stage architecture (verify effects, then check authorization)---but it uses an oracle, not a real verifier. The constructive question is whether a non-oracle verifier can approach this bound.

\subsection{Constructive Framework: EffectVerif-AuthMonitor}

We evaluate the EffectVerif-AuthMonitor framework (Section~3.4) by replacing the oracle present label with deterministic execution-trace and static verifiers, then measuring whether the authorization monitor can detect unauthorized effects on held-out trace groups.

\begin{table}[t]
\centering
\caption{EffectVerif-AuthMonitor results. Thresholds selected on validation trace groups, evaluated on held-out trace groups. Static verifier uses tool-call rules without trace evidence. These are controlled diagnostics, not deployment calibration.}
\label{tab:verifier}
\resizebox{\textwidth}{!}{
\begin{tabular}{llcccc}
\toprule
Setting & Trace family & Rows & Exec verifier FNR/FPR & Static verifier FNR/FPR \\
\midrule
T58/T62 & Controlled auth traces & 1,104 & 0.0000 / 0.0000 & 0.0000 / 0.0520 \\
T59/T62 & real-agent-tools local adapters & 384 & 0.1765 / 0.0000 & 0.1765 / 0.0000 \\
T61/T62 & DeepSeek provider API & 960 & 0.0000 / 0.0000 & 1.0000 / 0.0000 \\
T63 & Broader local adapters & 2,400 & 0.0351 / 0.0000 & 0.6140 / 0.0132 \\
T65 & Headless Chrome runtime & 960 & 0.0000 / 0.0000 & 1.0000 / 0.0000 \\
T64 & Live HTTP + webhook protocol & 2,400 & 0.0000 / 0.0000 & 0.6294 / 0.0462 \\
\bottomrule
\end{tabular}}
\end{table}

Table~\ref{tab:verifier} shows a consistent pattern: execution-trace verifiers substantially outperform static verifiers, especially on provider, broader-tool, browser-runtime, and protocol effects. On the live protocol traces, the execution verifier reaches test FNR/FPR of $0.0/0.0$ ($N^+=197$, $N^-=1170$) while the static verifier has FNR/FPR of $0.629/0.046$. On the headless Chrome browser-runtime traces, execution verification reaches $0.0/0.0$ ($N^+=62$, $N^-=411$) while static rules miss all browser effects (FNR $1.0$). Static verifiers fail because they rely on tool-name-to-effect mappings; when the same effect is produced through a tool whose name does not trigger the rule, or when the effect occurs remotely (provider API), the static rule has no signal.

\textbf{Honest limits.} Four important caveats apply. (1) \textbf{Verifier-generator alignment.} The execution verifier is deterministic and rule-based, and was designed with knowledge of the controlled trace formats used in this study. The FNR/FPR of $0.0/0.0$ on T64/T65 is partly a consequence of the verifier and trace generator sharing the same effect-to-evidence mapping assumptions. It should be interpreted as feasibility evidence that execution-level evidence can reduce unauthorized-effect misses \emph{when} verifier design and trace format are aligned---not as a performance estimate for unseen trace formats or deployed settings. Future work should replace this handcrafted verifier with an independently learned or formally specified one. (2) The live protocol traces cover real outbound HTTPS and local webhook I/O, but not provider-backed search, SaaS messaging, or HTTP browser automation. The browser-runtime traces use actual headless Chrome but over file-backed DOM/JavaScript pages, not live web browsing. (3) As Layer~3 showed, row-level success does not translate to action-level safety, and simple calibration fails. (4) Several LOTO cells have small positive counts ($N^+ = 8, 10, 20$); Wilson 95\% CIs are $\pm 0.10$--$0.25$ for these cells. The highest FNR values (e.g., $0.97$ for \texttt{file\_written}) have wider CIs and should be treated as stress-test indicators rather than precise point estimates.

\textbf{Contrastive projection as observed-pair upper bound.} Full-training contrastive projection, which uses cross-form positive pairs from all tools, achieves post-FNR of $0.00$ on 5/6 effects and reduces the worst case (\texttt{content\_fetched}) from $0.78$ to $0.25$. A strict leave-one-pair-out protocol improves 40/68 tool-case evaluations (mean $\Delta$FNR $=+0.157$). These results confirm that the fragmentation is at least partially representation-geometric and can be repaired when cross-form pairs are observed. However, since strict train-only projection does not beat strong baselines (Section~4.5), we use contrastive projection only as evidence that representation alignment can repair observed surface pairs. Full contrastive tables, multiseed results, and alignment theory are in the appendix.

\subsection{Existing-Defense Proxy Comparison}

Table~\ref{tab:proxies} compares EffectVerif-AuthMonitor against proxy monitors that approximate existing defense paradigms without faithfully reimplementing them.

\begin{table}[t]
\centering
\caption{Proxy defense comparison (mean across all held-out trace groups). Pre-action rule-only denies based on tool name and policy rules without trace evidence. Provenance-only checks whether the call context supports the action. Raw-status boundary checks exit codes and explicit error status strings.}
\label{tab:proxies}
\begin{tabular}{lcccc}
\toprule
Monitor & Mean row FNR & Mean row FPR & Mean U-Allow & Mean FDeny \\
\midrule
Pre-action rule-only & 0.5700 & 0.0186 & 0.5014 & 0.1627 \\
Provenance-only & 0.2873 & 0.0733 & 0.0486 & 0.3762 \\
EffectVerif (label-hidden) & 0.1162 & 0.0480 & 0.0459 & 0.7428 \\
Raw-status boundary & 0.1113 & 0.0101 & 0.0171 & 0.0265 \\
\bottomrule
\end{tabular}}
\end{table}

Pre-action rule-only and provenance-only monitors miss many realized-effect violations (FNR $0.57$ and $0.29$). EffectVerif with label-hidden evidence improves substantially on these. However, a simple handcrafted raw-status boundary---which checks only exit codes and explicit error status strings---is competitive (row FNR $0.1113$, U-Allow $0.0171$, FDeny $0.0265$). This competitiveness is concentrated on effects with clear status signatures: \texttt{tool\_error} maps directly to non-zero exit codes, and \texttt{command\_executed} is trivially detectable. Raw-status boundary cannot detect effects without status signatures---it would miss \texttt{network\_egress} via a successful curl call, \texttt{content\_fetched} via a 200-OK API response, or \texttt{file\_written} via a zero-exit write. EffectVerif's advantage is thus effect-dependent and should be characterized per-effect, not aggregated. The contribution is the Auth-SafeInv evaluation framework and the three-layer diagnosis, not a universal claim of method superiority over status-policy defenses.

\section{Discussion}

Our results paint a consistent picture: agent safety monitoring evaluation can be systematically over-optimistic, and the over-optimism operates at distinct layers that require distinct diagnostic methods. We draw four broader implications.

\textbf{Auth-SafeInv makes a specific failure mode visible.} When safety is evaluated as attack success rate, task completion, or pooled F1, unauthorized-effect misses under coverage gaps can be invisible. Auth-SafeInv surfaces this failure mode by conditioning the target on task authorization and requiring held-out surface-form evaluation. The Layer~1 evidence is a SafeInv diagnostic (can the representation support effect-invariant readout?) rather than a full Auth-SafeInv claim (does the monitor detect unauthorized effects under task authorization?). We maintain this distinction because LOTO fragmentation could also be influenced by domain shift (different argument formats, context templates, and effect co-occurrence patterns across tools), not only by representation-level effect fragmentation. The convergent evidence---lexical controls, pIIA cross-form transfer patterns, and cross-model consistency---partially constrains this alternative, but does not eliminate it. We therefore treat Layer~1 as an operational diagnostic: a monitor that fails held-out-tool SafeInv cannot satisfy the stronger Auth-SafeInv condition.

\textbf{Each layer can mask the next.} The three layers we identify are not independent. A monitor that appears safe at Layer~1 (low within-tool FNR) can fail at Layer~2 (the low FNR depends on structured labels). A monitor that passes Layer~2 (reasonable label-hidden FNR) can fail at Layer~3 (row-level metrics do not translate to safe action policies). In our controlled setting, row-level FNR/FPR of $0.0/0.0$ on provider and browser traces coexists with action-level FDeny of $1.0$ and all-allow calibration collapse. We do not claim that all agent safety evaluations exhibit these failures---our evidence is limited to the protocols and monitors we constructed. We recommend that agent safety monitor evaluations in general report at minimum: (a) held-out tool or trace-family FNR/FPR, (b) a label-field ablation comparing full and minimal trace views, and (c) action-level U-Allow and FDeny alongside row-level metrics.

\textbf{Negative results are evidence.} The three gate failures (T54/T55/T56) are not null results---they are evidence that pure representation-side repair, at least in the forms we tested, is insufficient for this problem. Reporting them prevents the literature from retrying approaches that fail under strict evaluation. We encourage the field to report such gate failures explicitly, with clear protocols and strong baselines, so that method development targets the right difficulty level.

\textbf{Constructive frameworks need honest limits.} EffectVerif-AuthMonitor is not a deployed safety solution. It works under controlled settings, depends on structured execution evidence, and has unresolved action-level calibration. But it serves as a constructive baseline: a concrete instantiation of the two-stage architecture (verify effects, then check authorization) against which future methods can be compared. The framework's limits---static verifier failures on provider/browser effects, competitive raw-status baselines, all-allow calibration collapse---also define the research agenda.

\section{Limitations}

\begin{enumerate}
    \item \textbf{Data.} The original diagnostic uses synthetic scenarios; lexical artifacts remain a risk despite controls. Auth-SafeInv traces are controlled (local sandbox, local adapters, single-provider API, file-backed browser, key-free live protocol) and do not cover provider-backed search, SaaS messaging, HTTP browser automation, or deployed-agent runtime logs.
    \item \textbf{Model scale.} All experiments use models $\leq 8$B parameters. Larger-model behavior is unknown.
    \item \textbf{Surface-form coverage.} Our fragmentation diagnosis is restricted mainly to tool switching ($\mathcal{S} = \mathcal{T}$). Extensions to context injection, semantic reframing, and language variation require separate validation.
    \item \textbf{Domain shift confounding.} LOTO splits change tool name, argument format, context template, and effect co-occurrence patterns simultaneously with the tool surface form. The observed FNR degradation therefore cannot be attributed exclusively to representation-level effect fragmentation. Lexical controls, pIIA, and cross-model consistency provide convergent but not decisive evidence.
    \item \textbf{Verifier-generator alignment.} The execution verifier is deterministic, rule-based, and designed for the controlled trace formats used in this study. The low FNR/FPR on T64/T65 partly reflects generator-verifier schema alignment. Independent learned or formally specified effect verifiers are needed before deployment-facing claims.
    \item \textbf{Action-level policy.} We show that simple threshold calibration fails (T73), but do not provide a working solution. Abstention policies, utility-aware calibration, and human-in-the-loop designs remain future work.
    \item \textbf{External defense comparison.} Our proxy baselines are conceptual approximations, not faithful reimplementations of external systems. The raw-status boundary baseline is competitive on effects with clear status signatures, but cannot detect effects without them.
    \item \textbf{Statistical.} Several key FNR estimates have small positive counts (e.g., $N^+=8$ for \texttt{content\_fetched/terminal}) and wide confidence intervals ($\pm 0.17$ for that cell). We report Wilson CIs and recommend treating small-$N$ cells as stress-test indicators rather than precise estimates.
    \item \textbf{Two-phase model assumption.} Our framework assumes tool calls execute in a reversible sandbox before commit. If sandbox and production environments produce different execution evidence, the verifier's decisions may not transfer.
\end{enumerate}

\section{Reproducibility}

All experiments are implemented in Python with PyTorch and scikit-learn. The original diagnostic pipeline is automated via \texttt{run\_experiments.sh}. The Auth-SafeInv pipeline is tracked by \texttt{reproduce\_auth\_safeinv.sh} and a result-source appendix mapping each reported artifact to its script, command, output files, and caveats. API-key-dependent provider traces are recorded as artifacts but not rerun by default. Code, data, and analysis outputs will be released as an open-source repository upon publication.

\bibliographystyle{plain}
\begin{thebibliography}{99}

\bibitem{jaw2026}
\textit{JAW: Comment and Control---Hijacking Agentic Workflows via Context-Grounded Evolution}. arXiv:2605.11229, May 2026.

\bibitem{litmus2026}
\textit{LITMUS: Benchmarking Behavioral Jailbreaks of LLM Agents in Real OS Environments}. arXiv:2605.10779, May 2026.

\bibitem{attrguard2026}
\textit{AttriGuard: Defeating Indirect Prompt Injection via Causal Attribution of Tool Invocations}. arXiv:2603.10749, Mar 2026.

\bibitem{clawguard2026}
\textit{ClawGuard: A Runtime Security Framework for Tool-Augmented LLM Agents}. arXiv:2604.11790, Apr 2026.

\bibitem{argus2026}
\textit{ARGUS: Defending LLM Agents Against Context-Aware Prompt Injection}. arXiv:2605.03378, May 2026.

\bibitem{xoa2026}
\textit{XOA: Execute-Only Agents}. Agentic OS Workshop, 2026.

\bibitem{pact2026}
\textit{PACT: The Granularity Mismatch in Agent Security: Argument-Level Provenance Solves Enforcement and Isolates the LLM Reasoning Bottleneck}. arXiv:2605.11039, May 2026.

\bibitem{alignmentcontracts2026}
\textit{Alignment Contracts for Agentic Security Systems}. arXiv:2605.00081, Apr 2026.

\bibitem{arm2026}
\textit{ARM: Causality Laundering---Denial-Feedback Leakage in Tool-Calling LLM Agents}. arXiv:2604.04035, Apr 2026.

\bibitem{agenttrust2026}
\textit{AgentTrust: Runtime Safety Evaluation and Interception for AI Agent Tool Use}. arXiv:2605.04785, May 2026.

\bibitem{parallax2026}
\textit{Parallax: Why AI Agents That Think Must Never Act}. arXiv:2604.12986, Apr 2026.

\bibitem{eca2026}
\textit{ECA: Hallucination as Exploit---Evidence-Carrying Multimodal Agents}. arXiv:2605.19192, May 2026.

\bibitem{lasm2026}
\textit{A Systematic Survey of Security Threats and Defenses in LLM-Based AI Agents}. arXiv:2604.23338, Apr 2026.

\bibitem{sok2026}
\textit{SoK: The Attack Surface of Agentic AI}. arXiv:2603.22928, Mar 2026.

\bibitem{queryipi2026}
\textit{QueryIPI: Query-agnostic Indirect Prompt Injection on Coding Agents}. arXiv:2510.23675, Jan 2026.

\bibitem{agentdojo2024}
\textit{AgentDojo: A Dynamic Environment to Evaluate Prompt Injection Attacks and Defenses for LLM Agents}. NeurIPS Datasets and Benchmarks, 2024.

\bibitem{toolemu2024}
\textit{Identifying the Risks of LM Agents with an LM-Emulated Sandbox}. ICLR, 2024.

\bibitem{agentharm2025}
\textit{AgentHarm: A Benchmark for Measuring Harmfulness of LLM Agents}. ICLR, 2025.

\bibitem{asb2025}
\textit{Agent Security Bench (ASB): Formalizing and Benchmarking Attacks and Defenses in LLM-based Agents}. ICLR, 2025.

\bibitem{stwebagentbench2026}
\textit{ST-WebAgentBench: A Benchmark for Evaluating Safety and Trustworthiness in Web Agents}. ICLR, 2026.

\bibitem{tracesafe2026}
\textit{TraceSafe: A Systematic Assessment of LLM Guardrails on Multi-Step Tool-Calling Trajectories}. arXiv:2604.07223, Apr 2026.

\bibitem{causalarmor2026}
\textit{CausalArmor: Efficient Indirect Prompt Injection Guardrails via Causal Attribution}. arXiv:2602.07918, Feb 2026.

\bibitem{irm2019}
M. Arjovsky, L. Bottou, I. Gulrajani, and D. Lopez-Paz. \textit{Invariant Risk Minimization}. arXiv:1907.02893, 2019.

\bibitem{groupdro2020}
S. Sagawa, P. W. Koh, T. B. Hashimoto, and P. Liang. \textit{Distributionally Robust Neural Networks for Group Shifts}. ICLR, 2020.

\bibitem{geiger2021}
A. Geiger et al. \textit{Causal Abstractions of Neural Networks}. NeurIPS, 2021.

\bibitem{geiger2022}
A. Geiger et al. \textit{Inducing Causal Structure for Interpretable Neural Networks}. ICML, 2022.

\bibitem{geiger2025}
A. Geiger et al. \textit{Causal Abstraction: A Theoretical Foundation for Mechanistic Interpretability}. JMLR, 2025.

\bibitem{sutter2025}
T. Sutter et al. \textit{The Non-Linear Representation Dilemma}. NeurIPS, 2025.

\bibitem{icidg2026}
J. Zhou et al. \textit{From Small to Large: In-Context Learning for Domain Generalization}. IJCV, 2026.

\bibitem{cdas2026}
Y. Bao et al. \textit{Faithful Bi-Directional Model Steering via Distribution Matching and DII}. ICLR, 2026.

\end{thebibliography}

\end{document}
